\chapter{Principles of Functional Analysis}
\lecture{25}{May 28.}{}

This chapter presents the most important general principles in functional analysis. These results explain why completeness is such a powerful assumption in the study of normed vector spaces. The first principle is the Uniform Boundedness Principle, also known as the Banach--Steinhaus Theorem. It says that a pointwise bounded family of bounded linear operators on a Banach space is automatically uniformly bounded in operator norm. The second principle is the Open Mapping Theorem. It says that a surjective bounded linear operator between Banach spaces maps open sets to open sets. Its most important consequence is the Inverse Operator Theorem: a bijective bounded linear operator between Banach spaces has a bounded inverse. An equivalent formulation is the Closed Graph Theorem, which says that a linear operator between Banach spaces is bounded if and only if its graph is closed. The third principle is the Hahn--Banach Theorem. It says that bounded linear functionals defined on a subspace of a normed vector space can be extended to the whole space without increasing their norm. This theorem is the main tool for producing linear functionals and separating points from closed subspaces.

\section{The Baire Category Theorem}
The Baire Category Theorem is the topological foundation behind the Uniform Boundedness Principle and the Open Mapping Theorem.

\begin{definition}[Nowhere dense]
    Let \(X\) be a metric space. A subset \(E \subseteq X\) is called nowhere dense if
    \[
        \mathrm{int}\!\left(\overline{E}\right) = \varnothing.
    \]
\end{definition}

\begin{theorem}[Baire Category Theorem] \label{thm: baire category theorem}
    Let \(X\) be a nonempty complete metric space. If
    \[
        X = \bigcup_{n=1}^{\infty} F_n,
    \]
    where each \(F_n \subseteq X\) is closed, then at least one of the sets \(F_n\) has nonempty interior. Equivalently, a nonempty complete metric space cannot be written as a countable union of closed nowhere dense subsets.
\end{theorem}

\begin{proof}
    We prove the contrapositive. Suppose that each \(F_n\) is closed and has empty interior. We shall prove that
    \[
        X \neq \bigcup_{n=1}^{\infty} F_n.
    \]
    In other words, we will construct a point \(x \in X\) such that \(x \notin F_n\) for every \(n \in \mathbb{N}\).

    The proof is based on the following idea. Since \(F_1\) has empty interior, it cannot contain any open ball. Hence we can find a small closed ball avoiding \(F_1\). Inside that ball, since \(F_2\) also has empty interior, we can find a smaller closed ball avoiding \(F_2\). Continuing in this way, we construct a nested sequence of closed balls whose radii tend to zero. Completeness will then guarantee that these balls contain a limiting point. By construction, that limiting point avoids every \(F_n\).

    We begin the construction. Since \(F_1\) has empty interior, \(\mathrm{int}(F_1) = \varnothing\). Thus \(F_1\) contains no nonempty open ball. Since \(X\) is nonempty and \(X\) itself is open in the metric topology, it follows that \(F_1 \neq X\). Hence there exists a point \(y_1 \in X \setminus F_1\). Because \(F_1\) is closed, its complement \(X \setminus F_1\) is open. Therefore there exists \(\omega_1 > 0\) such that \(B_{\omega_1}(y_1) \subseteq X \setminus F_1\). Choose \(r_1 > 0\) so small that \(0 < r_1 < 1\) and \(r_1 < \omega_1\). Then
    \[
        \overline{B_{r_1}}(y_1) \subseteq B_{\omega_1}(y_1) \subseteq X \setminus F_1.
    \]
    Renaming \(y_1\) as \(x_1\), we have found a closed ball
    \[
        \overline{B_{r_1}}(x_1) \subseteq X \setminus F_1, \quad 0 < r_1 < 1.
    \]

    We now proceed inductively. Suppose that, for some \(n \geq 1\), we have already chosen a closed ball \(\overline{B_{r_n}}(x_n)\). Since \(F_{n+1}\) has empty interior, it cannot contain the nonempty open ball \(B_{r_n}(x_n)\). Therefore there exists a point \(y_{n+1} \in B_{r_n}(x_n) \setminus F_{n+1}\). Because \(F_{n+1}\) is closed, the complement \(X \setminus F_{n+1}\) is open. Hence there exists \(\omega_{n+1} > 0\) such that \(B_{\omega_{n+1}}(y_{n+1}) \subseteq X \setminus F_{n+1}\). Also, since \(y_{n+1} \in B_{r_n}(x_n)\), the open ball \(B_{r_n}(x_n)\) is a neighborhood of \(y_{n+1}\). Thus, by choosing the radius sufficiently small, we may ensure that the next closed ball lies inside both \(B_{r_n}(x_n)\) and \(X \setminus F_{n+1}\). More precisely, choose \(r_{n+1} > 0\) such that
    \[
        0 < r_{n+1} < \frac{1}{n+1}
    \]
    and \(\overline{B_{r_{n+1}}}(y_{n+1}) \subseteq B_{r_n}(x_n) \cap (X \setminus F_{n+1})\). Renaming \(y_{n+1}\) as \(x_{n+1}\), we obtain
    \[
        \overline{B_{r_{n+1}}}(x_{n+1}) \subseteq B_{r_n}(x_n) \setminus F_{n+1},
    \]
    with
    \begin{equation} \label{eq: baire radius bound}
        0 < r_{n+1} < \frac{1}{n+1}.
    \end{equation}
    By induction, we obtain a sequence of closed balls \(\overline{B_{r_1}}(x_1), \overline{B_{r_2}}(x_2), \overline{B_{r_3}}(x_3), \ldots\) such that
    \begin{equation} \label{eq: baire ball nesting}
        \overline{B_{r_{n+1}}}(x_{n+1}) \subseteq \overline{B_{r_n}}(x_n) \quad \text{for every } n \in \mathbb{N},
    \end{equation}
    and
    \begin{equation} \label{eq: baire ball disjoint}
        \overline{B_{r_n}}(x_n) \cap F_n = \varnothing \quad \text{for every } n \in \mathbb{N}.
    \end{equation}
    Moreover, by \autoref{eq: baire radius bound}, \(r_n \to 0\).

    We claim that \((x_n)\) is a Cauchy sequence. Let \(m \geq n\). Since the balls are nested, \autoref{eq: baire ball nesting} implies that \(x_m \in \overline{B_{r_n}}(x_n)\). Therefore \(d(x_m, x_n) \leq r_n\). Since \(r_n \to 0\), it follows that for every \(\varepsilon > 0\), there exists \(N \in \mathbb{N}\) such that \(r_N < \varepsilon\). Hence, whenever \(m \geq n \geq N\),
    \[
        d(x_m, x_n) \leq r_n \leq r_N < \varepsilon.
    \]
    Thus \((x_n)\) is Cauchy. Since \(X\) is complete, there exists \(x \in X\) such that \(x_n \to x\).

    We now show that this limit point avoids every \(F_n\). Fix \(n \in \mathbb{N}\). For every \(m \geq n\), the nestedness condition gives \(x_m \in \overline{B_{r_n}}(x_n)\). Since \(\overline{B_{r_n}}(x_n)\) is closed and \(x_m \to x\), we may pass to the limit and obtain \(x \in \overline{B_{r_n}}(x_n)\). But by \autoref{eq: baire ball disjoint}, the ball \(\overline{B_{r_n}}(x_n)\) is disjoint from \(F_n\). Hence \(x \notin F_n\). Since \(n \in \mathbb{N}\) was arbitrary, we conclude that
    \[
        x \notin \bigcup_{n=1}^{\infty} F_n.
    \]
    Therefore \(X \neq \bigcup_{n=1}^{\infty} F_n\). This proves the contrapositive. Consequently, if a nonempty complete metric space \(X\) is written as a countable union of closed sets, \(X = \bigcup_{n=1}^{\infty} F_n\), then at least one of the sets \(F_n\) must have nonempty interior.
\end{proof}

\section{The Uniform Boundedness Principle}
Let \(X\) be a set. A family \(\{f_i\}_{i \in I}\) of functions \(f_i : X \to Y_i\), where each \(Y_i\) is a normed vector space, is called \textbf{pointwise bounded} if
\[
    \sup_{i \in I} \lVert f_i(x) \rVert_{Y_i} < \infty \quad \text{for every } x \in X.
\]

\begin{theorem}[Uniform Boundedness Principle] \label{thm: uniform boundedness principle}
    Let \(X\) be a Banach space, let \(I\) be an index set, and, for each \(i \in I\), let \(Y_i\) be a normed vector space. Suppose \(A_i : X \to Y_i\) is a bounded linear operator for every \(i \in I\). If the family \(\{A_i\}_{i \in I}\) is pointwise bounded, that is,
    \[
        \sup_{i \in I} \lVert A_i x \rVert_{Y_i} < \infty \quad \text{for every } x \in X,
    \]
    then
    \[
        \sup_{i \in I} \lVert A_i \rVert < \infty.
    \]
\end{theorem}

\begin{proof}
    For \(n \in \mathbb{N}\), define
    \[
        F_n \coloneqq \left\{ x \in X : \sup_{i \in I} \lVert A_i x \rVert_{Y_i} \leq n \right\}.
    \]
    Each \(F_n\) is closed. Indeed, if \(x_m \in F_n\) and \(x_m \to x\), then for each \(i \in I\),
    \[
        \lVert A_i x \rVert_{Y_i} = \lim_{m \to \infty} \lVert A_i x_m \rVert_{Y_i} \leq n,
    \]
    and hence \(x \in F_n\). By pointwise boundedness,
    \[
        X = \bigcup_{n=1}^{\infty} F_n.
    \]
    Since \(X\) is complete, the Baire Category Theorem implies that some \(F_N\) has nonempty interior. Hence there exist \(x_0 \in X\) and \(\varepsilon > 0\) such that \(B_\varepsilon(x_0) \subseteq F_N\). Thus
    \[
        \lVert A_i x \rVert_{Y_i} \leq N \quad \text{for every } i \in I \text{ and every } x \in B_\varepsilon(x_0).
    \]
    We now show that this local uniform bound implies a uniform bound on operator norms. Let \(x \in X\) with \(x \neq 0\). Then both points
    \[
        x_0 + \frac{\varepsilon}{2} \frac{x}{\lVert x \rVert} \quad \text{and} \quad x_0 - \frac{\varepsilon}{2} \frac{x}{\lVert x \rVert}
    \]
    belong to \(B_\varepsilon(x_0)\). Hence, by the local uniform bound, for every \(i \in I\),
    \[
        \left\lVert A_i \!\left( x_0 + \frac{\varepsilon}{2} \frac{x}{\lVert x \rVert} \right) \right\rVert_{Y_i} \leq N, \qquad \left\lVert A_i \!\left( x_0 - \frac{\varepsilon}{2} \frac{x}{\lVert x \rVert} \right) \right\rVert_{Y_i} \leq N.
    \]
    Using the linearity of \(A_i\), we have
    \[
        A_i \!\left( x_0 + \frac{\varepsilon}{2} \frac{x}{\lVert x \rVert} \right) - A_i \!\left( x_0 - \frac{\varepsilon}{2} \frac{x}{\lVert x \rVert} \right) = A_i \!\left( \frac{\varepsilon x}{\lVert x \rVert} \right) = \frac{\varepsilon}{\lVert x \rVert} A_i x.
    \]
    Therefore, by the triangle inequality,
    \[
        \frac{\varepsilon}{\lVert x \rVert} \lVert A_i x \rVert_{Y_i} \leq 2N.
    \]
    Equivalently, \(\lVert A_i x \rVert_{Y_i} \leq \frac{2N}{\varepsilon} \lVert x \rVert\). Thus \(\lVert A_i \rVert \leq \frac{2N}{\varepsilon}\) for every \(i \in I\). Consequently, \(\sup_{i \in I} \lVert A_i \rVert < \infty\).
\end{proof}

\begin{corollary}[Pointwise limits of bounded operators] \label{cl: pointwise limits of bounded operators}
    Let \(X\) be a Banach space and let \(Y\) be a normed vector space. Suppose \(A_n : X \to Y\) is a sequence of bounded linear operators such that, for every \(x \in X\), the limit
    \[
        Ax \coloneqq \lim_{n \to \infty} A_n x
    \]
    exists in \(Y\). Then \(A : X \to Y\) is a bounded linear operator, and
    \[
        \lVert A \rVert \leq \liminf_{n \to \infty} \lVert A_n \rVert.
    \]
\end{corollary}

\begin{proof}
    For each \(x \in X\), the sequence \((A_n x)\) converges, and hence is bounded. Therefore
    \[
        \sup_{n \in \mathbb{N}} \lVert A_n x \rVert_Y < \infty
    \]
    for every \(x \in X\). By the Uniform Boundedness Principle, \(\sup_{n \in \mathbb{N}} \lVert A_n \rVert < \infty\). The map \(A\) is linear because it is defined as a pointwise limit of linear maps. Moreover, for every \(x \in X\),
    \[
        \lVert Ax \rVert_Y = \lim_{n \to \infty} \lVert A_n x \rVert_Y = \liminf_{n \to \infty} \lVert A_n x \rVert_Y \leq \left( \liminf_{n \to \infty} \lVert A_n \rVert \right) \lVert x \rVert_X.
    \]
    Taking the supremum over \(\lVert x \rVert_X \leq 1\), we obtain \(\lVert A \rVert \leq \liminf_{n \to \infty} \lVert A_n \rVert\).
\end{proof}

\section{The Open Mapping Theorem}
A map \(T : X \to Y\) between topological spaces is called \textbf{open} if \(T(U)\) is open in \(Y\) whenever \(U\) is open in \(X\).

\begin{theorem}[Open Mapping Theorem] \label{thm: open mapping theorem}
    Let \(X\) and \(Y\) be Banach spaces. If \(T : X \to Y\) is a surjective bounded linear operator, then \(T\) is open.
\end{theorem}

We prove the theorem in two steps.

\begin{lemma} \label{lm: open mapping step 1}
    Let \(X\) and \(Y\) be Banach spaces, and let \(T : X \to Y\) be a surjective bounded linear operator. Let \(B_X \coloneqq \{x \in X : \lVert x \rVert < 1\}\). Then there exists \(\delta > 0\) such that
    \[
        B_Y(0, \delta) \subseteq \overline{T(B_X)}.
    \]
\end{lemma}

\begin{proof}
    Since \(T\) is surjective, we have \(Y = T(X)\). Also,
    \[
        X = \bigcup_{n=1}^{\infty} n B_X,
    \]
    because for every \(x \in X\), one can choose \(n > \lVert x \rVert_X\), and then \(x \in n B_X\). Hence
    \[
        Y = T(X) = \bigcup_{n=1}^{\infty} T(n B_X) = \bigcup_{n=1}^{\infty} n T(B_X),
    \]
    where the last equality follows from the linearity of \(T\). Therefore
    \[
        Y = \bigcup_{n=1}^{\infty} \overline{n T(B_X)}.
    \]
    Indeed, since \(n T(B_X) \subseteq \overline{n T(B_X)}\), the inclusion \(Y \subseteq \bigcup_{n=1}^{\infty} \overline{n T(B_X)}\) follows, while the reverse inclusion holds because each \(\overline{n T(B_X)}\) is a subset of \(Y\). Each set \( \overline{n T(B_X)}\) is closed in \(Y\). Since \(Y\) is a Banach space, it is a complete metric space. Hence, by the Baire Category Theorem, there exists \(n_0 \in \mathbb{N}\) such that
    \[
        \mathrm{int}\!\left( \overline{n_0 T(B_X)} \right) \neq \varnothing.
    \]
    Since scalar multiplication by \(n_0\) is a homeomorphism of \(Y\), we have \( \overline{n_0 T(B_X)} = n_0 \overline{T(B_X)}\). Thus \(\overline{T(B_X)}\) also has nonempty interior.

    Consequently, there exist \(y_0 \in Y\) and \(r > 0\) such that \(B_Y(y_0, r) \subseteq \overline{T(B_X)}\). We now show that this implies that \(\overline{T(B_X)}\) contains an open ball centered at the origin. Since \(B_X\) is convex and symmetric, the set \(T(B_X)\) is also convex and symmetric. Therefore its closure \(\overline{T(B_X)}\) is convex and symmetric.

    Let \(y \in Y\) satisfy \(\lVert y \rVert_Y < r\). Then
    \[
        y_0 + y \in B_Y(y_0, r) \quad \text{and} \quad y_0 - y \in B_Y(y_0, r).
    \]
    Hence both \(y_0 + y\) and \(y_0 - y\) lie in \(\overline{T(B_X)}\). By symmetry of \(\overline{T(B_X)}\), \(-(y_0 - y) = -y_0 + y \in \overline{T(B_X)}\). Since \(\overline{T(B_X)}\) is convex,
    \[
        y = \frac{1}{2}(y_0 + y) + \frac{1}{2}(-y_0 + y) \in \overline{T(B_X)}.
    \]
    Thus \(B_Y(0, r) \subseteq \overline{T(B_X)}\). Taking \(\delta = r\) proves the lemma.
\end{proof}

\begin{lemma} \label{lm: open mapping step 2}
    Let \(X\) and \(Y\) be Banach spaces, and let \(T : X \to Y\) be a bounded linear operator. Suppose that for some \(\delta > 0\),
    \[
        B_Y(0, \delta) \subseteq \overline{T(B_X)}.
    \]
    Then
    \[
        B_Y(0, \delta) \subseteq T(B_X).
    \]
\end{lemma}

\begin{proof}
    Let \(y \in Y\) satisfy \(\lVert y \rVert_Y < \delta\). Choose a number \(\sigma\) such that \(\lVert y \rVert_Y < \sigma < \delta\). We shall construct a sequence \((x_n)_{n=1}^{\infty}\) in \(X\) such that
    \[
        \sum_{n=1}^{\infty} \lVert x_n \rVert_X < 1 \quad \text{and} \quad y = \sum_{n=1}^{\infty} T x_n.
    \]
    Then \(x \coloneqq \sum_{n=1}^{\infty} x_n\) will belong to \(B_X\), and by continuity of \(T\), we will have \(Tx = y\).

    We use the assumption \(B_Y(0, \delta) \subseteq \overline{T(B_X)}\) in the following scaled form. If \(z \in Y\) and \(a > 0\) satisfy \(\lVert z \rVert_Y < a\delta\), then \(\lVert z/a \rVert_Y < \delta\), so \(z/a \in \overline{T(B_X)}\). Therefore, for every \(\varepsilon > 0\), there exists \(u \in B_X\) such that \(\lVert z/a - Tu \rVert_Y < \varepsilon/a\). Setting \(x \coloneqq au\), we get \(\lVert x \rVert_X < a\) and
    \[
        \lVert z - Tx \rVert_Y = a\lVert z/a - Tu \rVert_Y < \varepsilon.
    \]
    Thus we have proved the following useful fact:
    \[
        \lVert z \rVert_Y < a\delta \implies \text{for every } \varepsilon > 0,\ \exists\, x \in X \text{ s.t. } \lVert x \rVert_X < a \text{ and } \lVert z - Tx \rVert_Y < \varepsilon.
    \]

    We now begin the construction. Since \(\lVert y \rVert_Y < \sigma = \delta \cdot \frac{\sigma}{\delta}\), we apply the above with \(z = y\), \(a = \frac{\sigma}{\delta}\), and \(\varepsilon = \frac{\delta - \sigma}{2}\). Thus there exists \(x_1 \in X\) such that
    \[
        \lVert x_1 \rVert_X < \frac{\sigma}{\delta} \quad \text{and} \quad \lVert y - T x_1 \rVert_Y < \frac{\delta - \sigma}{2}.
    \]
    Define \(r_1 \coloneqq y - T x_1\). Then \(\lVert r_1 \rVert_Y < \frac{\delta - \sigma}{2}\).

    We now carry out the induction. Suppose that, for some \(n \geq 1\), we have already chosen \(x_1, \ldots, x_n \in X\) such that the remainder
    \[
        r_n \coloneqq y - \sum_{k=1}^n T x_k
    \]
    satisfies \(\lVert r_n \rVert_Y < \frac{\delta - \sigma}{2^n}\). Set \(a_n \coloneqq \frac{\delta - \sigma}{\delta \cdot 2^n}\). Then \(a_n \delta = \frac{\delta - \sigma}{2^n}\), and by the induction hypothesis, \(\lVert r_n \rVert_Y < a_n \delta\). Therefore we may apply the above with \(z = r_n\), \(a = a_n\), and \(\varepsilon = \frac{\delta - \sigma}{2^{n+1}}\). Hence there exists \(x_{n+1} \in X\) such that
    \[
        \lVert x_{n+1} \rVert_X < \frac{\delta - \sigma}{\delta \cdot 2^n} \quad \text{and} \quad \left\lVert y - \sum_{k=1}^{n+1} T x_k \right\rVert_Y < \frac{\delta - \sigma}{2^{n+1}}.
    \]
    This completes the induction.

    We now estimate the series \(\sum x_n\). Using the above bounds,
    \begin{align*}
        \sum_{n=1}^{\infty} \lVert x_n \rVert_X = \lVert x_1 \rVert_X + \sum_{n=1}^{\infty} \lVert x_{n+1} \rVert_X < \frac{\sigma}{\delta} + \sum_{n=1}^{\infty} \frac{\delta - \sigma}{\delta \cdot 2^n} = \frac{\sigma}{\delta} + \frac{\delta - \sigma}{\delta} = 1.
    \end{align*}
    Thus the series \(\sum_{n=1}^{\infty} x_n\) is absolutely convergent in \(X\). Since \(X\) is Banach, it converges to some element \(x \coloneqq \sum_{n=1}^{\infty} x_n \in X\). Moreover, \(\lVert x \rVert_X \leq \sum_{n=1}^{\infty} \lVert x_n \rVert_X < 1\), hence \(x \in B_X\). Since \(T\) is bounded, it is continuous. Therefore
    \[
        Tx = T\!\left( \sum_{n=1}^{\infty} x_n \right) = \sum_{n=1}^{\infty} T x_n.
    \]
    On the other hand, the remainder estimate gives \(\left\lVert y - \sum_{k=1}^n T x_k \right\rVert_Y < \frac{\delta - \sigma}{2^n} \to 0\), hence \(\sum_{n=1}^{\infty} T x_n = y\). Consequently, \(Tx = y\), and since \(x \in B_X\), we have \(y \in T(B_X)\). This proves \(B_Y(0, \delta) \subseteq T(B_X)\).
\end{proof}

Now we can prove open mapping theorem.
\begin{proof}[Proof of \autoref{thm: open mapping theorem}]
    By \autoref{lm: open mapping step 1}, there exists \(\delta > 0\) such that \(B_Y(0, \delta) \subseteq \overline{T(B_X)}\). By \autoref{lm: open mapping step 2}, \(B_Y(0, \delta) \subseteq T(B_X)\).

    Let \(U \subseteq X\) be open. We prove that \(T(U)\) is open. Choose \(y_0 \in T(U)\), say \(y_0 = Tx_0\) with \(x_0 \in U\). Since \(U\) is open, there exists \(\varepsilon > 0\) such that \(B_X(x_0, \varepsilon) \subseteq U\). We claim that \(B_Y(y_0, \delta\varepsilon) \subseteq T(U)\). Indeed, if \(\lVert y - y_0 \rVert_Y < \delta\varepsilon\), then \(\lVert \varepsilon^{-1}(y - y_0) \rVert_Y < \delta\). Hence there exists \(u \in X\) such that \(\lVert u \rVert_X < 1\) and \(Tu = \varepsilon^{-1}(y - y_0)\). Therefore \(y = y_0 + \varepsilon Tu = T(x_0 + \varepsilon u)\). Since \(x_0 + \varepsilon u \in B_X(x_0, \varepsilon) \subseteq U\), we get \(y \in T(U)\). Hence \(T(U)\) is open.
\end{proof}

\section{The Inverse Operator and Closed Graph Theorems}

\begin{theorem}[Inverse Operator Theorem] \label{thm: inverse operator theorem}
    Let \(X\) and \(Y\) be Banach spaces. If \(T : X \to Y\) is a bijective bounded linear operator, then \(T^{-1} : Y \to X\) is bounded.
\end{theorem}

\begin{proof}
    Since \(T\) is surjective and bounded, the Open Mapping Theorem implies that \(T\) is open. Therefore, for every open set \(U \subseteq X\),
    \[
        (T^{-1})^{-1}(U) = T(U)
    \]
    is open in \(Y\). Hence \(T^{-1}\) is continuous. Since \(T^{-1}\) is linear, it is bounded.
\end{proof}

\begin{definition}[Graph of a linear operator]
    Let \(X\) and \(Y\) be normed vector spaces, and let \(T : X \to Y\) be linear. The graph of \(T\) is
    \[
        \mathrm{graph}(T) \coloneqq \left\{ (x, Tx) : x \in X \right\} \subseteq X \times Y.
    \]
    We equip \(X \times Y\) with the norm
    \[
        \lVert (x, y) \rVert_{X \times Y} \coloneqq \lVert x \rVert_X + \lVert y \rVert_Y.
    \]
    If \(X\) and \(Y\) are Banach spaces, then \(X \times Y\) is Banach.
\end{definition}

\begin{theorem}[Closed Graph Theorem] \label{thm: closed graph theorem}
    Let \(X\) and \(Y\) be Banach spaces. Let \(T : X \to Y\) be a linear operator. Then \(T\) is bounded if and only if \(\mathrm{graph}(T)\) is closed in \(X \times Y\).
\end{theorem}

\begin{proof}
    First suppose that \(T\) is bounded. Let \((x_n, Tx_n) \to (x, y)\) in \(X \times Y\). Then \(x_n \to x\) in \(X\), and \(Tx_n \to y\) in \(Y\). Since \(T\) is continuous, \(Tx_n \to Tx\). Hence \(y = Tx\), and so \((x, y) \in \mathrm{graph}(T)\). Thus the graph is closed.

    Conversely, suppose that \(\mathrm{graph}(T)\) is closed. Since \(X \times Y\) is Banach, the closed subspace \(\mathrm{graph}(T)\) is also Banach. Consider the projection
    \[
        \pi_X : \mathrm{graph}(T) \to X, \quad \pi_X(x, Tx) = x.
    \]
    This map is linear, bijective, and bounded because
    \[
        \lVert \pi_X(x, Tx) \rVert_X = \lVert x \rVert_X \leq \lVert x \rVert_X + \lVert Tx \rVert_Y.
    \]
    By the Inverse Operator Theorem, \(\pi_X^{-1}\) is bounded. Hence there exists \(C > 0\) such that
    \[
        \lVert x \rVert_X + \lVert Tx \rVert_Y   = \lVert (x, Tx) \rVert_{X \times Y} = \lVert \pi_X^{-1} (x) \rVert_{X \times Y} \leq C \lVert x \rVert_X.
    \]
    Therefore \(\lVert Tx \rVert_Y \leq (C - 1) \lVert x \rVert_X\) for every \(x \in X\). Thus \(T\) is bounded.
\end{proof}

\begin{remark}
    The Closed Graph Theorem is often used in the following form: to prove that a linear map \(T : X \to Y\) between Banach spaces is bounded, it suffices to prove that
    \[
        x_n \to x, \quad Tx_n \to y \implies Tx = y.
    \]
\end{remark}

\section{The Hahn--Banach Theorem}
We now turn to the extension theorem for bounded linear functionals.

\begin{definition}[Sublinear function]
    Let \(X\) be a real vector space. A function \(p : X \to \mathbb{R}\) is called sublinear if
    \[
        p(x + y) \leq p(x) + p(y) \quad \text{for all } x, y \in X,
    \]
    and
    \[
        p(\lambda x) = \lambda p(x) \quad \text{for all } x \in X,\ \lambda \geq 0.
    \]
\end{definition}

\begin{theorem}[Hahn--Banach theorem, real form] \label{thm: hahn-banach real form}
    Let \(X\) be a real vector space, let \(M \subseteq X\) be a linear subspace, and let \(p : X \to \mathbb{R}\) be sublinear. Suppose \(f : M \to \mathbb{R}\) is linear and satisfies
    \[
        f(x) \leq p(x) \quad \text{for every } x \in M.
    \]
    Then there exists a linear functional \(F : X \to \mathbb{R}\) such that
    \[
        F\vert_M = f \quad \text{and} \quad F(x) \leq p(x) \quad \text{for every } x \in X.
    \]
\end{theorem}

\begin{proof}
    We first prove the one-dimensional extension step. Let \(z \in X \setminus M\) and set \(M_1 \coloneqq M + \mathbb{R} z\). We seek an extension of the form
    \[
        f_1(x + tz) = f(x) + t\alpha, \quad x \in M,\ t \in \mathbb{R}.
    \]
    We need \(f(x) + t\alpha \leq p(x + tz)\) for all \(x \in M\), \(t \in \mathbb{R}\). For \(t > 0\), this is equivalent to \(\alpha \leq p(u + z) - f(u)\) for all \(u \in M\). For \(t < 0\), it is equivalent to \(\alpha \geq f(u) - p(u - z)\) for all \(u \in M\). Thus it suffices to choose \(\alpha\) such that
    \[
        f(u) - p(u - z) \leq \alpha \leq p(v + z) - f(v) \quad \text{for all } u, v \in M.
    \]
    Such an \(\alpha\) exists because
    \[
        f(u) + f(v) = f(u + v) \leq p(u + v) \leq p(u - z) + p(v + z),
    \]
    and hence \(f(u) - p(u - z) \leq p(v + z) - f(v)\). Therefore we may choose
    \[
        \alpha \in \left[ \sup_{u \in M}\!\left( f(u) - p(u - z) \right),\ \inf_{v \in M}\!\left( p(v + z) - f(v) \right) \right].
    \]
    With this choice of \(\alpha\), we have, for every \(x \in M\),
    \begin{equation} \label{eq: hb one-dim pos}
        f(x) + \alpha \leq p(x + z),
    \end{equation}
    and
    \begin{equation} \label{eq: hb one-dim neg}
        f(x) - \alpha \leq p(x - z).
    \end{equation}
    We now explain why these two inequalities, which correspond to the coefficients \(1\) and \(-1\) of the new direction \(z\), imply the desired estimate for all real coefficients \(t\). The important point is that \autoref{eq: hb one-dim pos} and \autoref{eq: hb one-dim neg} hold for every element of \(M\), not just for one fixed element.

    Let \(x \in M\) and \(t \in \mathbb{R}\). If \(t > 0\), then \(x + tz = t(x/t + z)\) and \(x/t \in M\). Applying \autoref{eq: hb one-dim pos} to \(x/t\) and multiplying by \(t > 0\), using the linearity of \(f\) and the positive homogeneity of \(p\), gives \(f(x) + t\alpha \leq p(x + tz)\). If \(t < 0\), write \(t = -s\) where \(s > 0\). Then \(x + tz = s(x/s - z)\) and \(x/s \in M\). Applying \autoref{eq: hb one-dim neg} to \(x/s\) and multiplying by \(s > 0\), using the linearity of \(f\) and the positive homogeneity of \(p\), gives \(f(x) - s\alpha \leq p(x - sz)\). Since \(t = -s\), this is exactly \(f(x) + t\alpha \leq p(x + tz)\). If \(t = 0\), the desired inequality reduces to \(f(x) \leq p(x)\), which holds by assumption on \(M\). Thus
    \[
        f_1(x + tz) = f(x) + t\alpha \leq p(x + tz) \quad \text{for all } x \in M,\ t \in \mathbb{R}.
    \]
    Therefore \(f_1\) extends \(f\) and satisfies \(f_1 \leq p\) on \(M_1\).

    Now let \(\mathcal{E}\) be the collection of all pairs \((N, g)\), where \(M \subseteq N \subseteq X\) is a linear subspace and \(g : N \to \mathbb{R}\) is a linear functional such that \(g\vert_M = f\) and \(g(x) \leq p(x)\) for every \(x \in N\). Thus \(\mathcal{E}\) consists of all possible extensions of \(f\) to larger subspaces, still dominated by \(p\). We define a partial order on \(\mathcal{E}\) by extension: for \((N_1, g_1), (N_2, g_2) \in \mathcal{E}\), we write \((N_1, g_1) \leq (N_2, g_2)\) if \(N_1 \subseteq N_2\) and \(g_2\vert_{N_1} = g_1\).

    Notice first that \(\mathcal{E}\) is nonempty, because \((M, f) \in \mathcal{E}\).

    We claim that every chain in \(\mathcal{E}\) has an upper bound. Let \(\mathcal{C} = \{(N_i, g_i)\}_{i \in I}\) be a chain in \(\mathcal{E}\). This means that for any \(i, j \in I\), either \((N_i, g_i) \leq (N_j, g_j)\) or \((N_j, g_j) \leq (N_i, g_i)\). Define \(N_{\mathcal{C}} \coloneqq \bigcup_{i \in I} N_i\). We first check that \(N_{\mathcal{C}}\) is a linear subspace of \(X\). Let \(x, y \in N_{\mathcal{C}}\) and \(a, b \in \mathbb{R}\). Then there exist \(i, j \in I\) such that \(x \in N_i\) and \(y \in N_j\). Since \(\mathcal{C}\) is a chain, the two subspaces \(N_i\) and \(N_j\) are comparable. For example, suppose \(N_i \subseteq N_j\). Then \(x, y \in N_j\), and since \(N_j\) is a linear subspace, \(ax + by \in N_j \subseteq N_{\mathcal{C}}\). The other case \(N_j \subseteq N_i\) is identical. Hence \(N_{\mathcal{C}}\) is a linear subspace.

    We now define a functional \(g_{\mathcal{C}} : N_{\mathcal{C}} \to \mathbb{R}\) as follows. Given \(x \in N_{\mathcal{C}}\), choose \(i \in I\) such that \(x \in N_i\), and set \(g_{\mathcal{C}}(x) \coloneqq g_i(x)\). This definition is well-defined. Indeed, if \(x \in N_i \cap N_j\), then because \(\mathcal{C}\) is a chain, the pairs \((N_i, g_i)\) and \((N_j, g_j)\) are comparable. Suppose, for instance, that \((N_i, g_i) \leq (N_j, g_j)\). Then \(g_j\vert_{N_i} = g_i\), and therefore \(g_j(x) = g_i(x)\). The other case is the same. Hence \(g_{\mathcal{C}}(x)\) is independent of the choice of \(i\).

    Next we verify that \(g_{\mathcal{C}}\) is linear. Let \(x, y \in N_{\mathcal{C}}\) and \(a, b \in \mathbb{R}\). Choose \(i, j \in I\) such that \(x \in N_i\) and \(y \in N_j\). Since \(\mathcal{C}\) is a chain, after possibly interchanging \(i\) and \(j\), we may assume that \((N_i, g_i) \leq (N_j, g_j)\). Then \(x, y \in N_j\), and hence \(ax + by \in N_j\). Therefore
    \[
        g_{\mathcal{C}}(ax + by) = g_j(ax + by) = a g_j(x) + b g_j(y) = a g_{\mathcal{C}}(x) + b g_{\mathcal{C}}(y).
    \]
    Thus \(g_{\mathcal{C}}\) is linear. Moreover, \(g_{\mathcal{C}}\) extends \(f\), since every \((N_i, g_i) \in \mathcal{E}\) satisfies \(M \subseteq N_i\) and \(g_i\vert_M = f\). Finally, \(g_{\mathcal{C}}(x) = g_i(x) \leq p(x)\) for any \(x \in N_i\). Therefore \((N_{\mathcal{C}}, g_{\mathcal{C}}) \in \mathcal{E}\). It is also an upper bound for \(\mathcal{C}\), because for every \(i \in I\), \(N_i \subseteq N_{\mathcal{C}}\) and \(g_{\mathcal{C}}\vert_{N_i} = g_i\).

    Thus every chain in \(\mathcal{E}\) has an upper bound. By Zorn's lemma, there exists a maximal element \((N_0, g_0) \in \mathcal{E}\). We claim that \(N_0 = X\). Suppose, for contradiction, that \(N_0 \neq X\). Then there exists \(z \in X \setminus N_0\). Since \(z \notin N_0\), the subspace \(N_0 + \mathbb{R} z\) is strictly larger than \(N_0\). By the one-dimensional extension step, there exists a linear functional \(\widetilde{g} : N_0 + \mathbb{R} z \to \mathbb{R}\) such that \(\widetilde{g}\vert_{N_0} = g_0\) and \(\widetilde{g}(x) \leq p(x)\) for every \(x \in N_0 + \mathbb{R} z\). Therefore \((N_0 + \mathbb{R} z, \widetilde{g}) \in \mathcal{E}\). Moreover, \((N_0, g_0) < (N_0 + \mathbb{R} z, \widetilde{g})\), and this extension is strict because \(z \notin N_0\). This contradicts the maximality of \((N_0, g_0)\).

    Hence \(N_0 = X\). Consequently, \(g_0 : X \to \mathbb{R}\) is a linear extension of \(f\) to all of \(X\), and it satisfies \(g_0(x) \leq p(x)\) for every \(x \in X\). This is the desired Hahn--Banach extension.
\end{proof}

\begin{theorem}[Hahn--Banach Extension Theorem] \label{thm: hahn-banach extension}
    Let \(X\) be a real normed vector space, and let \(M \subseteq X\) be a linear subspace. If \(f : M \to \mathbb{R}\) is a bounded linear functional, then there exists \(F \in X^*\) such that
    \[
        F\vert_M = f \quad \text{and} \quad \lVert F \rVert_{X^*} = \lVert f \rVert_{M^*}.
    \]
\end{theorem}

\begin{proof}
    Define \(p(x) \coloneqq \lVert f \rVert_{M^*} \lVert x \rVert\) for \(x \in X\). Then \(p\) is sublinear. For \(x \in M\),
    \[
        f(x) \leq \vert f(x) \vert \leq \lVert f \rVert_{M^*} \lVert x \rVert = p(x).
    \]
    By the Hahn--Banach theorem, there exists a linear functional \(F : X \to \mathbb{R}\) such that \(F\vert_M = f\) and \(F(x) \leq p(x)\) for every \(x \in X\). Applying this inequality to \(-x\), we also get \(-F(x) \leq p(x)\). Therefore
    \[
        \vert F(x) \vert \leq \lVert f \rVert_{M^*} \lVert x \rVert \quad \text{for every } x \in X.
    \]
    Thus \(F\) is bounded and \(\lVert F \rVert_{X^*} \leq \lVert f \rVert_{M^*}\). The reverse inequality follows from \(F\vert_M = f\). Hence \(\lVert F \rVert_{X^*} = \lVert f \rVert_{M^*}\).
\end{proof}

\begin{remark}[Complex Hahn--Banach theorem]
    The complex version is also true. If \(X\) is a complex normed vector space and \(M \subseteq X\) is a complex linear subspace, then every bounded complex-linear functional \(f : M \to \mathbb{C}\) extends to a bounded complex-linear functional \(F : X \to \mathbb{C}\) with
    \[
        F\vert_M = f, \quad \lVert F \rVert = \lVert f \rVert.
    \]
    One proves this by applying the real Hahn--Banach theorem to \(\mathrm{Re}\, f\) and then reconstructing the complex-linear functional.
\end{remark}

\section{Basic Consequences of Hahn--Banach}

\begin{corollary}[Existence of norming functionals] \label{cl: existence of norming functionals}
    Let \(X\) be a normed vector space and let \(x_0 \in X\), \(x_0 \neq 0\). Then there exists \(F \in X^*\) such that
    \[
        \lVert F \rVert = 1 \quad \text{and} \quad F(x_0) = \lVert x_0 \rVert.
    \]
\end{corollary}

\begin{proof}
    Let \(M = \mathrm{span}\{x_0\}\), and define \(f : M \to \mathbb{R}\) by \(f(tx_0) \coloneqq t \lVert x_0 \rVert\). Then
    \[
        \vert f(tx_0) \vert = \vert t \vert \lVert x_0 \rVert = \lVert tx_0 \rVert,
    \]
    so \(\lVert f \rVert = 1\). By the Hahn--Banach Extension Theorem, \(f\) extends to some \(F \in X^*\) with \(\lVert F \rVert = 1\). Since \(F(x_0) = f(x_0) = \lVert x_0 \rVert\), the result follows.
\end{proof}

\begin{corollary}[Dual norm formula] \label{cl: dual norm formula}
    Let \(X\) be a normed vector space. Then for every \(x \in X\),
    \[
        \lVert x \rVert = \sup\!\left\{ \vert F(x) \vert : F \in X^*,\ \lVert F \rVert \leq 1 \right\}.
    \]
\end{corollary}

\begin{proof}
    For every \(F \in X^*\) with \(\lVert F \rVert \leq 1\), \(\vert F(x) \vert \leq \lVert x \rVert\). Thus
    \[
        \sup\!\left\{ \vert F(x) \vert : \lVert F \rVert \leq 1 \right\} \leq \lVert x \rVert.
    \]
    If \(x = 0\), equality is clear. If \(x \neq 0\), \autoref{cl: existence of norming functionals} gives \(F \in X^*\) such that \(\lVert F \rVert = 1\) and \(F(x) = \lVert x \rVert\). Hence equality holds.
\end{proof}

\begin{corollary}[Separation from a closed subspace] \label{cl: separation from closed subspace}
    Let \(X\) be a normed vector space, let \(M \subseteq X\) be a closed linear subspace, and let \(x_0 \notin M\). Then there exists \(F \in X^*\) such that
    \[
        F\vert_M = 0 \quad \text{and} \quad F(x_0) \neq 0.
    \]
    More precisely, one can choose \(F \in X^*\) such that
    \[
        \lVert F \rVert = 1, \quad F(x_0) = \mathrm{dist}(x_0, M).
    \]
\end{corollary}

\begin{proof}
    Let \(d \coloneqq \mathrm{dist}(x_0, M) = \inf_{m \in M} \lVert x_0 - m \rVert\). Since \(M\) is closed and \(x_0 \notin M\), we have \(d > 0\). Let \(N \coloneqq M + \mathbb{R} x_0\). Define \(f : N \to \mathbb{R}\) by
    \[
        f(m + tx_0) \coloneqq td, \quad m \in M,\ t \in \mathbb{R}.
    \]
    This is well-defined and linear. Also, \(f\vert_M = 0\) and \(f(x_0) = d\). We claim that \(\lVert f \rVert \leq 1\). If \(t \neq 0\), then
    \[
        \lVert m + tx_0 \rVert = \vert t \vert \left\lVert x_0 + \frac{m}{t} \right\rVert \geq \vert t \vert d = \vert f(m + tx_0) \vert.
    \]
    The case \(t = 0\) is immediate. Thus \(\lVert f \rVert \leq 1\). By the Hahn--Banach Extension Theorem, extend \(f\) to \(F \in X^*\) with \(\lVert F \rVert = \lVert f \rVert \leq 1\). Since \(F(x_0) = d\) and \(F\) vanishes on \(M\), we have for every \(m \in M\),
    \[
        d = \vert F(x_0 - m) \vert \leq \lVert F \rVert \lVert x_0 - m \rVert.
    \]
    Taking the infimum over \(m \in M\), we get \(d \leq \lVert F \rVert \cdot d\). Since \(d > 0\), this implies \(\lVert F \rVert \geq 1\). Therefore \(\lVert F \rVert = 1\).
\end{proof}

\paragraph*{Summary.}
The essential principles of this chapter are:
\begin{itemize}
    \item \textbf{The Baire Category Theorem:} completeness prevents a space from being too small in a countable-union sense.
    \item \textbf{The Uniform Boundedness Principle:} pointwise bounded families of bounded linear operators on a Banach space are uniformly bounded in operator norm.
    \item \textbf{The Open Mapping Theorem:} surjective bounded linear operators between Banach spaces are open.
    \item \textbf{The Inverse Operator Theorem:} bijective bounded linear operators between Banach spaces have bounded inverses.
    \item \textbf{The Closed Graph Theorem:} a linear operator between Banach spaces is bounded if and only if its graph is closed.
    \item \textbf{The Hahn--Banach Theorem:} bounded linear functionals extend from subspaces to the whole space without increasing their norm.
\end{itemize}
Together, these results form the basic toolkit for working with Banach spaces. They show that completeness converts algebraic hypotheses into strong continuity and extension properties.
